Differentiable pose-estimation layers must differentiate a globally optimal pose, which is discontinuous at branch crossings. Deterministic layers built on semidefinite programming (SDP) are certifiable when the relaxation is tight, but they lift the problem to a matrix variable and have no defined behaviour at a crossing. Boltzmann-distribution layers avoid hard branch selection but usually rely on Monte Carlo integration in six dimensions. The quaternion formulation of quadratic pose-estimation problems (QPEP) supplies the algebraic structure for a third option. For the rotation subproblem left after exact translation elimination, we prove that the standard $10$-dimensional rotation lift is \emph{exactly} $2Uv_2(q)$ for an orthogonal $U$, that the 21 $SO(3)$ constraints span precisely the degree-two vanishing ideal of the Veronese variety $\nu_2(\mathbb{RP}^3)$, and that the resulting $10\times10$ SDP is the Veronese (Shor) relaxation of a quartic on $\mathbb{S}^3$; a constructed Choi--Lam instance has a certified gap of $0.034188$ with $\mathrm{rank}(X^\star)=6$. We derive a closed-form Euler multiplier $\lambda^\star=4J$ and a $3\times3$ Riemannian implicit differentiation on $\mathbb{S}^3$ with exact antipodal invariance. Replacing the hard argmin by a multi-modal Laplace posterior---translation integrated out exactly, rotation approximated analytically over \emph{all} QPEP minima with weights $w_k\propto e^{-J_k/\tau}/\sqrt{\det H_k}$---gives a layer with no Monte Carlo variance, no SDP lift, and one $3\times3$ solve per active mode. On 15 real LINEMOD frames the hard argmin attains a median $2.06^\circ$ and $3.7$ mm ADD-S, while the Laplace posterior mean is pulled toward the competing mode ($44.24^\circ$, $13.9$ mm, paired bias $+36.9^\circ$); a total-degree homotopy confirms the generic $40$ complex projective roots on $4/4$ instances.
